\section{Discussion}
\label{sec:discussion}

The main conclusion is conditional but sharp within its regular experiment. In a finite-query teacher experiment, a functionally identifiable TT quotient pays approximation error plus gauge-adjusted dimension over queries. A horizontal lower packing shows that the dimension term is not a proof artifact. Under the stated local information-normal and off-chart separation conditions, approximation and estimation are simultaneously unavoidable for a fixed compression class. The local constant depends on the relation between query and deployment information; \(d/2\) is the matched conditional-KL instance of the classical local asymptotic minimax calculation.

The access model can change the answer. White-box access to a feasible teacher gives zero risk immediately; exact soft labels also give finite-sample recovery in the identifiable logistic construction of \cref{prop:access}. These examples refute a universal positive \(d/n\) lower bound after the observation channel is changed, but they do not characterize the minimax rate of every soft-label experiment. Under misspecification, random deployment prompts create a quadrature term indexed by \(\tr(V_XH^{-1})\), which can be far smaller than \(d\) or zero. This distinction should influence practice: query-limited API distillation, finite Monte Carlo teacher sampling, exact soft-label distillation, and white-box tensorization should not automatically share one sample-complexity claim. Statistical explanations of soft-label variance reduction already appear in the distillation literature~\citep{phuong2019distillation,menon2021distillation,xie2026distillation}; our contribution is the access warning within the TT query experiment.

The gauge calculation is also insufficient by itself. It removes redundant core coordinates, but only restricted functional curvature makes all remaining tensor directions statistically visible. The results here cover the nondegenerate case in which the functional information is uniformly positive on the horizontal tangent. They do not prove a \(d_{\mathrm{fun}}/n\) law merely by replacing dimension with the rank of a singular pullback; weak eigenvalues and higher-order identification can produce different rates. Developing localized spectral complexity and adaptive penalties for weakly excited quotient directions is a natural next theoretical problem.

The rank-boundary result separates two goals that are often conflated. Penalized likelihood may adapt in prediction risk without consistently identifying a minimal rank. Honest rank identification requires functional separation from lower strata and, for consistency, a penalty stronger than an AIC-scale constant. The present two-point argument is generic; a TT-specific strengthening would relate functional separation quantitatively to the smallest relevant TT-unfolding singular values. A general theory over singular unions of TT varieties, rather than separated manifolds, would be required for a broader extension.

Same-data capability certification is useful only when the whole fitting map is stable. Its band is deliberately conservative and is an application of classical stability and bounded-differences arguments, not a new concentration principle~\citep{bousquet2002stability,feldman2019stability,bousquet2020sharper}. PAC--Bayes and protected-test methods already provide alternative simultaneous multi-risk guarantees~\citep{adams2024multiple,angelopoulos2021ltt,laufer2023pareto}. Cross-fitting, selective likelihood, or cluster jackknife methods may improve the adaptation price, but a final full-data refit creates a target mismatch that must be resolved rather than ignored. Protected-holdout Learn-then-Test remains preferable when nonconvex optimization has no stability guarantee.

Our literature search through 7 September 2026 supports a narrow priority statement, not proof of universal absence: among the primary sources examined, we did not identify a result simultaneously treating common TT-compressed weights inside a frozen nonlinear conditional generator, population teacher-relative conditional KL across environments, tensorization/rank selection in a finite stochastic-query experiment, and a local lower bound indexed by information on the TT gauge quotient. Tree-tensor risk oracles, TT minimax and inference, nonlinear low-rank neural rates, population/minimax KV-cache compression, statistical distillation, and Fisher- or KL-aware rank allocation already occupy each broader claim separately~\citep{michel2022tree,zhou2022ttoi,qin2024ttregression,arunan2026lora,haverbeck2026kvcache,menon2021distillation,iosvd2026,laser2026}. Our priority claim is therefore restricted to this coupling and to its explicit observation-model and functional-identifiability boundaries, rather than to any ingredient alone.

\paragraph{Limitations.}
The bounded likelihood-ratio assumption excludes unrestricted long-sequence log loss; the nonlinear verification covers a frozen block rather than every transformer layer; the statistical result assumes a global or certified approximate empirical optimizer; and the experiments are controlled tangent models rather than evidence from large open models. In addition, the leading oracle theorem is restricted to compact rank-separated strata and its constants can deteriorate with the atlas and smallest TT unfolding singular values. The joint approximation--estimation lower result constructs only a local information-normal tube and requires a separated off-chart complement. These are substantive limitations and define the next proof and empirical tasks.

\paragraph{Conclusion.}
We have formulated population functional risk as an explicit statistical target for query-limited TT compression and have shown why dimension-based guarantees require gauge quotienting, functional visibility, a fixed teacher access model, and controlled quotient charts. Under the stated regular, local information-normal, and off-chart separation conditions, the modelwise approximation--estimation frontier is sharp; outside them, local Fisher surrogates and \(d/n\) heuristics can be misleading.
